\section{Exceptional ends: three discriminants and the exact intersections}
The finite input to this section is fully specified, rather than hidden
behind a group-name lookup. The program \texttt{exceptional\_amalgams.g}
uses exact GAP subgroup and finitely presented group algorithms. It first
enumerates every ordered generating pair $(a,b)$ in $A_4,S_4,A_5$ with
$|a|,|b|>2$, $|ab|=2$ and $\langle a\rangle\cap\langle b\rangle=1$.
The numbers of pairs are respectively $24,48,360$, and their types modulo
automorphisms are
\[
 A_4:(3,3),\qquad S_4:(3,4),(4,3),\qquad
 A_5:(3,5),(5,3),(5,5).
\]
For every $(5,5)$ pair in $A_5$ it checks
$a b^2a^3b^2a^2=1$. This is a complete finite premise, not a sample of
projective groups. The standard spherical triangle presentations identify
the $(3,3),(3,4),(3,5)$ groups as $A_4,S_4,A_5$.

\begin{lemma}\label{E.ExceptionalTypes}
If $p\equiv3\pmod4$ is prime and $L(p)$ is a chiral rank-four rotation
group, its type, up to duality, is $(5,3,4)$, $(5,3,5)$ or $(3,5,3)$.
\end{lemma}
\begin{proof}
Neither proper end fixes a point. A dihedral end with both rotation
orders greater than two is impossible: both rotations would lie in its
cyclic rotation subgroup, and their trivial cyclic intersection together
with $(ab)^2=1$ would force their orders to be at most two. The retained
prime-field subgroup theorem therefore leaves $A_4,S_4,A_5$.
The displayed complete local pair census leaves only the compatible
three-label types. The types $(3,3,3)$, $(3,3,4)$, $(3,4,3)$ and
$(3,3,5)$ and their duals are excluded by their spherical universal
rotation groups. These have respectively order $60,192,576,7200$;
the latter has only $C_2$ and $A_5$ composition factors, and the middle
two are solvable. No simple quotient is $L(p)$ with $p\equiv3\pmod4$.
These are the standard simplex, cube, $24$-cell and $120$-cell groups;
their full presentation orders and composition factors are also checked
by the same finite program.

For $(4,3,4)$ the trace quadratic has $a^2=c^2=2$, $b^2=1$ and
$\Delta=0$. Its explicit projective reversal excludes chirality.
The remaining unwanted types are $(5,5,3)$ and $(5,5,5)$ and their duals.
Add the preceding five-five relator to each five-five end of its full
rotation presentation. The exact finite calculation gives order $7200$
and a composition series with successive orders $7200,120,2,1$ in both
cases. Their only nonabelian simple composition factor is $A_5$.
Again $L(p)$ cannot be a quotient. This exhausts every compatible local
pair, without an unrecorded search in $L(p)$.
\end{proof}

For completeness we give the full finite intersection test. Write
\[
 T(r,s,t)=\langle a,b,c\mid a^r=b^s=c^t=(ab)^2=(bc)^2=(abc)^2=1\rangle.
\]
In an $A_5$ end of type $(5,3)$ or $(3,5)$ the word
$u(a,b)=ab^2a^2b^2a^2$ is an involution inverting $b$ and generates,
with $b$, its full dihedral normalizer. In an $S_4$ end of type $(3,4)$
the corresponding word is $v(c,b)=cb^2c^2$. All these word assertions
are checked on every relevant local pair. The only proper overgroups
of $\langle b\rangle$ in $A_5$ have orders $3,6,12$ when $|b|=3$,
and $5,10$ when $|b|=5$. Thus a dihedral collision between the two ends
forces $u(c,b)=u(a,b)b^k$ (using $v$ in the $S_4$ case). For
$k=1,\ldots,s$ the exact quotient orders are
\[
\begin{array}{c|c}
 (r,s,t)& |T(r,s,t)/\langle\!\langle\hbox{collision}\rangle\!\rangle|\\\hline
 (5,3,4)&60,1,1\\
 (5,3,5)&3420,60,60\\
 (3,5,3)&660,1,360,360,1.
\end{array}
\]
For an $A_4$ collision put
$u_1=a^2b^{-1}a$, $u_2=ab^{-1}a^2$.
The two subgroups $\langle u_i,b\rangle$ are exactly the $A_4$
overgroups in the first $A_5$. Their three involutions are the
$b$-conjugates of $u_ib$. At the other end use $v=b^2c^2$ for $S_4$,
or $v_1=c^2b^{-1}c$, $v_2=cb^{-1}c^2$ for $A_5$.
It is therefore enough, and is exhaustive, to add each relation
$u_i^{b^j}=v_l$, $j=0,1,2$.
The six $(5,3,4)$ quotients all have order two.
For $(5,3,5)$ the twelve orders, in lexicographic $(i,l,j)$ order, are
\[
 1,1,1,\quad1,60,60,\quad60,1,60,\quad1,1,1.
\]
This corrects the source's assertion in Lemma~5.12 that all twelve
quotients are trivial. Its exclusion of $A_4$ intersections outside
$L(4)\cong L(5)$ remains valid. Full finitely presented group orders,
not orders of selected permutation images, are used in this calculation.
The nontrivial collision groups are identified as $A_5$, $L(9)$, $L(11)$
and $L(19)$ at orders $60,360,660,3420$. Their simplicity shows that the
complete intersection conclusions hold for arbitrary projective quotients:
exclude $q=4,5$ for $(5,3,4)$; $q=4,5,19$ for $(5,3,5)$; and $q=9,11$
for $(3,5,3)$. Thus the source's three intersection lemmas, not merely
their prime-field applications, follow from this one finite interface.

\begin{theorem}\label{E.PrimeExceptional}
Let $p\ge7$ be prime. For each square root of $5$ in $F_p$ define
\[
 D_{534}=1\pm\sqrt5,\qquad
 D_{535}=(7\pm5\sqrt5)/2,\qquad
 D_{353}=3\pm2\sqrt5.
\]
For type $(5,3,4)$ also require that $2$ is a square in $F_p$.
For any of these three types, the number of ordered chiral polytope
isomorphism classes with rotation group $L(p)$ is twice the number
of nonzero squares in the corresponding two-element list.
In the $(3,5,3)$ case this formula gives zero at $p=11$; in the
$(5,3,5)$ case it gives two at $p=19$. A missing square root of $5$
means there are no such classes.
The same types with group $L(q)$, $q\equiv3\pmod4$, require $q=p$ prime.
\end{theorem}
\begin{proof}
A local $(5,3)$ or $(3,5)$ pair generates $A_5$ by the spherical
triangle presentation and simplicity. A $(3,4)$ pair generates $S_4$:
the exact order four prevents its proper triangle quotients. Fix
traces as above. The nonzero quantity $d=a^2+b^2-4$ follows either
from the absolute irreducibility of the $A_5$ end or directly from
the quadratic equation for its order-five trace. The trace extension
theorem gives all extensions and their conjugacy classes.
For $(5,3,4)$ it gives $D_{534}$, while for $(3,5,3)$ it gives $D_{353}$.
For $(5,3,5)$, if $a^2=c^2$, it gives $D_{535}$.
If the two order-five trace squares differ, their sum is three and
product one, so $\Delta=1$. Both resulting extensions lie inside the
first $A_5$. Indeed, for the fixed pair
$a=(1\,2\,3\,4\,5)$, $b=(2\,5\,3)$ the complete local enumeration
finds exactly two order-five elements $c$ satisfying both involution
relations, and both are in the other $A_5$ conjugacy class from $a$.
The two classes have different trace squares in every characteristic
under consideration: replacing an order-five element by its square
replaces its trace square $t$ by $(t-2)^2=3-t$, using
$t^2-3t+1=0$; these differ because the characteristic is not five. Pair uniqueness and the two-root formula show
that these are exactly the mixed-class extensions. They cannot generate
$L(p)$.

Every remaining extension generates $L(p)$. A subgroup containing
$A_5$ in a prime-field $L(p)$ is either that $A_5$ or the whole group,
by the retained subgroup classification. In type $(5,3,4)$, $c$ has
order four and hence is outside $A_5$. In type $(5,3,5)$ the preceding
complete local check excludes same-class $c$ inside $A_5$.
In type $(3,5,3)$ the complete local check finds no order-three $c$
inside the first $A_5$ satisfying the two involution relations.
Thus generation is uniform; no large-group enumeration is substituted
for it.

The intersection table proves the end intersection is exactly
$\langle b\rangle$, except possibly at $p=19$ in type $(5,3,5)$
and $p=11$ in type $(3,5,3)$. Adjacent cyclic intersections are trivial
by the coprime orders. The required program \texttt{exceptional\_small.g}
checks every root of the same equations at these two primes. At $p=19$
the two nonzero-root extensions have end intersection of order three;
the zero-root extension has intersection six. At $p=11$ the sole
$(3,5,3)$ extension has zero discriminant and intersection ten.
Thus every nonzero-root extension has the full rank-four rotation
intersection property, and none of the zero-root cases contributes.

For a prime field every automorphism of $L(p)$ is induced by a
projectivity, a retained classical input. The trace extension theorem
therefore rules out reversal for every nonzero root. It produces the
reversal at a zero root whether or not the intersection property holds.
The two choices of order-five trace square cannot be identified by
an automorphism, and the invariant $2\operatorname{tr}(AC)-ac$
distinguishes the two roots within a choice. Conversely pair uniqueness
shows that these are all classes. This proves the exact counts.

Finally suppose $q=p^d\equiv3\pmod4$ and such a type is realized.
Here $d$ is odd. Its local $A_5$ (and $S_4$ where applicable) traces
lie in an extension of the prime field of degree at most two, so they
already lie in $F_p$. The same applies to a square root of $\Delta$
inside $F_q$. Conjugate the first pair to the displayed prime-field
normal form. The unique extension formulas put the third matrix in
$\mathrm{SL}_2(p)$ as well. Hence the generated group is conjugate to
$L(p)$, and equality with $L(q)$ forces $d=1$.
\end{proof}

The residue conditions of the source now follow from ordinary quadratic
reciprocity. The products of the paired discriminants are respectively
$-4,-19,-11$. For $p\ne19,11$ the Legendre symbols of the latter two
are $(p/19)$ and $(p/11)$. The quadratic residues modulo $19$ are
$1,4,5,6,7,9,11,16,17$; those modulo $11$ are $1,3,4,5,9$.
If the product is a nonsquare exactly one discriminant is a square;
if it is a square both are squares or neither is.
For $(5,3,4)$ the order-five and order-four trace requirements give
$p\equiv1,9,31,39\pmod{40}$. At residues $31,39$ exactly one
$D_{534}$ is a square, and at $1,9$ either both or neither are.
These are the complete counts of source Theorems~6, 7, 9 and~11.
In particular residue $10\pmod{11}$ belongs to the nonsquare-product
case of Theorem~11. Its omission from Theorem~1(4b) in the exact
arXiv-v1 edition is a statement error, not an additional exception to
the construction. A separately recorded $p=571$ matrix witness makes
that omission explicit.

\begin{proposition}\label{E.AllPointedExistence}
For $q\ge4$, $P(q)$ occurs as the full group of a chiral rank-four
polytope exactly when $q>4$. For $q\equiv1\pmod4$, $L(q)$ occurs
exactly when $q>9$.
\end{proposition}
\begin{proof}
For $P(q)$ choose a primitive $t$, of order $q-1$. If $q>4$, its order
exceeds $p^{d-1}+1$ when $d>1$, so no nonidentity field automorphism
fixes or inverts it. The pointed theorem applies, with nonsquare $t$
in odd characteristic. For $L(q)$, $q\equiv1\pmod4$, choose $t$ of
order $(q-1)/2$. This is even and greater than two when $q>9$.
For $d>1$ it exceeds $p^{d-1}+1$ except at $p^d=9$: the required
strict inequality is $p^{d-1}(p-2)>3$. Again the pointed theorem
applies. At $d=1$ there is no field automorphism to test.

The omitted groups are $L(4)=P(4)=L(5)\cong A_5$ and $L(9)\cong A_6$.
The complete finite proof in \texttt{small\_rotation\_exclusions.g}
examines all pairs $(a,b)$ modulo the actual automorphism group, retaining
$|a|,|b|>2$, $|ab|=2$, trivial cyclic intersection and a proper first end.
For each it examines every $c=(ab)t$ with $t$ an involution; this is
exactly the condition $(abc)^2=1$. It then applies the remaining
relations, generation and complete rank-four intersection conditions.
There are one and seven pair representatives, respectively. The resulting
polytopal triples number one and four (before the final stabilizer
quotient), and each has a verified involutory automorphism
$(a,b,c)\mapsto(a^{-1},a^2b,c)$. Thus none is chiral. This is a complete
finite premise for these groups, separate from the uniform existence proof.
\end{proof}
